% ===================================================================
%  TABELUL Nr. 9 — Munții și stațiunea. — Pădurea și vânătoarea.
%  Traduction 2026 d'apres texto/io, controlee sur texto/fr, texto/en
%  et texto/es. Consigne : voir texto/ro/10-tabelul-01.tex.
%
%  LE TABLEAU DE LA MONTAGNE PROLONGE CELUI DE LA FERME, et pour la
%  meme raison : « munte », « stânca », « vale », « izvor », « pădure »,
%  « brad », « peștera » — le relief est latin ou tres ancien, comme le
%  troupeau, parce que c'est le meme monde et la meme continuite.
% ===================================================================

%%K t09-apar-1 apar
\VUsaut{4.0mm}
\VUtitre{15.0pt}{60}{TABELUL Nr. 9}
\VUsaut{3.0mm}

%%K t09-tit-1 sub
\VUcentre{12.0pt}{0}{\textit{Prima scenă.}}
\VUsaut{3.0mm}

%%K t09-tit-2 sub
\VUpk{11.4pt}{40}{Munții. --- Stațiunea balneară.}
\VUsaut{3.0mm}

%%K t09-c1-01-1 p
1. --- Sunt de o săptămână la Pratz. Nu pot spune toată uimirea pe care am
încercat-o stând înaintea minunatului \VUgras{lanț de munți}
\textsuperscript{(1)} al Pirineilor, a cărui \VUgras{creastă}
\textsuperscript{(2)} se taie pe cerul albastru ca un fierăstrău uriaș.

%%K t09-c1-02-1 p
2. --- Nu îmi închipuiam că \VUgras{vârfurile} pirineene \textsuperscript{(3)} pot
ajunge la o asemenea \VUgras{înălțime}, și am rămas uimit de \VUgras{ghețarii}
minunați \textsuperscript{(5)} și de \VUgras{piscurile} înzăpezite
\textsuperscript{(4)}, pe care se năpustesc \VUgras{avalanșe} atât de
înspăimântătoare.

%%K t09-tit-3 sub
\VUsaut{3.0mm}
\VUpk{10.2pt}{40}{Priveliște de munte.}
\VUsaut{3.0mm}

%%K t09-c1-03-1 p
3. --- Chiar a doua zi după sosire am mers la vechiul \VUgras{vulcan} stins
\textsuperscript{(6)} de lângă Pratz. Pe marginea goală și stearpă a
\VUgras{craterului} se văd încă limpede urmele lăsate de \VUgras{lavă}, care a curs
acum câteva veacuri pe \VUgras{coasta muntelui} \textsuperscript{(7)}. Pe una
dintre \VUgras{râpe} am izbutit să găsesc câteva bucăți din această lavă.

%%K t09-c1-04-1 p
4. --- Pratz stă la hotar, iar \VUgras{cetatea} \textsuperscript{(8)}, care
păzește \VUgras{trecătoarea} \textsuperscript{(27)} dintre Franța și Spania,
seamănă mult cu vechile castele de pe malurile Rinului, care parcă pândesc
\VUgras{prada} din vârful stâncii lor.

%%K t09-c1-05-1 p
5. --- Un \VUgras{vultur} uriaș \textsuperscript{(9)}, al cărui cuib este aproape
de \VUgras{fort} \textsuperscript{(8)}, îmi atrage adesea luarea-aminte. Această
\VUgras{pasăre răpitoare} cu \VUgras{cioc} puternic aduce adesea puilor un miel
sau un ied strâns în \VUgras{gheare}. \VUgras{Aripile} îi sunt atât de mari, încât
poate pluti multă vreme cu povara ei grea.

%%K t09-tit-4 sub
\VUsaut{3.0mm}
\VUpk{10.2pt}{40}{Drumețul.}
\VUsaut{3.0mm}

%%K t09-c1-06-1 p
6. --- Plimbarea cea mai plăcută de aici este drumul spre \VUgras{cascada}
\textsuperscript{(10)} de la Coo. \VUgras{Apa căzătoare} se prăvale în vârtej și
apoi se face nevăzută ca un \VUgras{pârâu} drăguț în \VUgras{pădurea de brazi}
\textsuperscript{(11)} de la apus de târg. Am urcat prin această pădure până la
\VUgras{stația meteorologică} \textsuperscript{(12)} și de pe vârful
\VUgras{turnului} am cuprins cu privirea toată \VUgras{panorama} ținutului.
\VUgras{Priveliștea} este minunată, și se pare că \VUgras{firea} a risipit asupra
acestui ținut toate comorile pe care le avea la îndemână.

%%K t09-c1-07-1 p
7. --- Ca să coborâm, am luat \VUgras{funicularul} \textsuperscript{(13)} și am
coborât la o \VUgras{stație} mică lângă \VUgras{ruinele} \textsuperscript{(14)}
vechiului \VUgras{castel feudal}, în care locuiau odinioară stăpânii Pratzului.

%%K t09-c1-08-1 p
8. --- Bietul castel! Ce trebuie să gândească văzând dedesubtul lui îngrijita
\VUgras{stațiune balneară} \textsuperscript{(15)}, ajunsă \VUgras{stațiune de ape}
la modă?

%%K t09-c1-08-2 p
\VUgras{Clima} este atât de plăcută aici, iar \VUgras{apa minerală} atât de
tămăduitoare, încât \VUgras{cura} care se face la \VUgras{sanatoriu}
\textsuperscript{(17)} este o adevărată plăcere.

%%K t09-tit-5 sub
\VUsaut{3.0mm}
\VUpk{10.2pt}{40}{Stațiune balneară plăcută.}
\VUsaut{3.0mm}

%%K t09-c1-09-1 p
9. --- Într-adevăr, s-a făcut totul ca șederea să fie plăcută. Au construit un
\VUgras{viaduct} mare \textsuperscript{(16)} pentru calea ferată;
\VUgras{parcul} \textsuperscript{(18)} de la \VUgras{cazinou}, unde se dau
\VUgras{concerte}, este întocmit fermecător. Câteva \VUgras{vile} elegante
\textsuperscript{(19)} au fost ridicate în jurul \VUgras{cazinoului}
\textsuperscript{(21)}, iar \VUgras{șoseaua} foarte bine ținută
\textsuperscript{(22)} face umblatul nespus de lesnicios.

%%K t09-c1-10-1 p
10. --- Un simplu \VUgras{terasament} \textsuperscript{(23)} desparte
\VUgras{drumul} \textsuperscript{(20)} de \VUgras{valea} însăși
\textsuperscript{(25)}, în fundul căreia se află un \VUgras{lăculeț} grațios
\textsuperscript{(26)}, care hrănește un \VUgras{pârâu} limpede
\textsuperscript{(24)}.

%%K t09-c1-11-1 p
11. --- De cealaltă parte a văii se lucrează o \VUgras{mină} de fier
\textsuperscript{(28)}. Am fost de câteva ori să văd cum coboară \VUgras{minerii}
în \VUgras{puț}, și am intrat chiar în \VUgras{galerie}, unde ei își petrec ziua
scoțând \VUgras{minereul}, în care se află \VUgras{metalul}.

%%K t09-c1-12-1 p
12. --- Lângă mină au construit o \VUgras{turnătorie} mare, ca să folosească pe
loc bogățiile scoase din ea. \VUgras{Furnalul} \textsuperscript{(29)}, încălzit cu
\VUgras{cărbunele} de aici din belșug, îngăduie să se scoată \VUgras{fonta}
repede și ieftin.

%%K t09-c1-13-1 p
13. --- Nu departe de mină se lucrează o \VUgras{carieră} \textsuperscript{(30)}.
Bătrânul \VUgras{cioplitor în piatră} desprinde cu \VUgras{târnăcopul} său nu
\VUgras{granit}, nu \VUgras{porfir} și nici măcar \VUgras{gresie}, ci simplă
\VUgras{piatră de construcție} pentru zidirea caselor.

%%K t09-c1-14-1 p
14. --- Una dintre cele mai frumoase podoabe ale acestui ținut este \VUgras{bradul}
\textsuperscript{(31)}. Pretutindeni --- în fundul \VUgras{râpei}
\textsuperscript{(32)}, pe malul \VUgras{șuvoiului} \textsuperscript{(33)}, lângă
\VUgras{prăpastie} \textsuperscript{(34)} sau lângă râpă --- acele lui subțiri își
aduc nota verde.

%%K t09-tit-6 sub
\VUsaut{3.0mm}
\VUpk{10.2pt}{40}{Plimbări pe jos.}
\VUsaut{3.0mm}

%%K t09-c1-15-1 p
15. --- Mă folosesc de vacanță ca să fac plimbări lungi. \VUgras{Potecile}
\textsuperscript{(35)} care șerpuiesc pe munte îngăduie să faci câțiva
\VUgras{kilometri}, iar adesea și câteva \VUgras{leghe}, fără să bagi de seamă
depărtarea.

%%K t09-c1-15-2 p
\VUgras{Bornele} \textsuperscript{(36)} și \VUgras{indicatoarele}
\textsuperscript{(37)} însemnează potecile, pe care întâlnești adesea un tânăr
\VUgras{muntean} \textsuperscript{(38)} cu \VUgras{desaga} \textsuperscript{(40)}
la șold, ducând o \VUgras{marmotă} \textsuperscript{(39)}, sau vreun
\VUgras{țigan} \textsuperscript{(41)}, la care \VUgras{căciula de blană}
\textsuperscript{(42)} se ceartă ciudat cu \VUgras{mantaua} veche și ruptă
\textsuperscript{(43)}. Duce în \VUgras{lanț} un \VUgras{urs} învățat
\textsuperscript{(44)}.

%%K t09-tit-7 sub
\VUpk{10.2pt}{40}{Ascensiuni.}
\VUsaut{3.0mm}

%%K t09-c1-16-1 p
16. --- Aproape în fiecare zi merg cu \VUgras{călăuza} \textsuperscript{(48)} la o
\VUgras{ascensiune}, și sunt foarte mândru când, cu \VUgras{ranița}
\textsuperscript{(47)} în spate și cu \VUgras{alpenstockul} în mână, ca un adevărat
\VUgras{turist} \textsuperscript{(46)}, iau cu asalt \VUgras{stâncile}
\textsuperscript{(45)}.

%%K t09-c1-17-1 p
17. --- De pe vârful muntelui oamenii de pe câmpie nu par mai mari decât furnicile;
\VUgras{turma de oi} \textsuperscript{(50)}, care paște pe \VUgras{pășune}
\textsuperscript{(49)}, seamănă cu o jucărie de copil. \VUgras{Pinul}
\textsuperscript{(51)} sau chiar \VUgras{cabana} de lemn \textsuperscript{(52)}
este un punct abia zărit pe verdeață.

%%K t09-c1-18-1 p
18. --- În aceste plimbări întâlnim adesea pe \VUgras{potecă}
\textsuperscript{(53)} \VUgras{vânătorul de capre negre} \textsuperscript{(54)},
purtând pe umăr animalul ucis adineauri.

%%K t09-c1-19-1 p
19. --- Am pus în ierbarul meu o frunză foarte deosebită de \VUgras{ferigă}
\textsuperscript{(55)} și o crenguță trandafirie drăguță de \VUgras{iarbă
neagră} \textsuperscript{(57)}, pe care am rupt-o la intrarea într-o \VUgras{peșteră}
mică \textsuperscript{(56)} la ultima plimbare.

%%K t09-tit-8 sub
\VUsaut{3.0mm}
\VUcentre{12.0pt}{0}{\textit{A doua scenă.}}
\VUsaut{3.0mm}

%%K t09-tit-9 sub
\VUpk{11.4pt}{40}{Pădurea. --- Vânătoarea.}
\VUsaut{3.0mm}

%%K t09-c2-01-1 p
1. --- Ieri am petrecut în pădure o zi încântătoare. Viața harnică a animalelor și
a \VUgras{insectelor} \textsuperscript{(5)} care trăiesc acolo mă interesa mult.

%%K t09-c2-01-2 p
\VUgras{Păianjenul} \textsuperscript{(1)}, țesându-și \VUgras{pânza}
\textsuperscript{(2)} între două crengi ca să prindă \VUgras{musca}
\textsuperscript{(3)} care trece în zbor; \VUgras{melcul} \textsuperscript{(4)},
târându-se greu cu căsuța lui; rădașca, pândind prada; furnica harnică, muncind
lângă \VUgras{furnicar} \textsuperscript{(6)} --- toate aceste vietăți mici mișunau
și trudeau cu o putere neobișnuită.

%%K t09-c2-02-1 p
2. --- \VUgras{Șarpele} \textsuperscript{(7)}, pe care mai întâi l-am luat drept
\VUgras{viperă}, iar el s-a dovedit un \VUgras{șarpe de casă} nevătămător, s-a
încolăcit sub trunchiul unui copac și din când în când își scotea căpșorul subțire,
mereu parcă gata să muște. Înaintea mea se târa greu un \VUgras{limax} mare
\textsuperscript{(8)}, care mâncase una dintre \VUgras{fragile} mici și gustoase
\textsuperscript{(11)}, de care sunt aici o mulțime.

%%K t09-c2-03-1 p
3. --- Pământul era așternut cu \VUgras{flori de pădure}: \VUgras{ciuboțica
cucului} \textsuperscript{(9)}, \VUgras{saschiul} \textsuperscript{(10)} și multe
alte plante necunoscute mie înfloreau printre \VUgras{smocurile de iarbă}
\textsuperscript{(12)}.

%%K t09-c2-04-1 p
4. --- În acest ținut \VUgras{trufa} este aproape tot atât de obișnuită ca
\VUgras{ciuperca} \textsuperscript{(14)}, iar \VUgras{purcelul}
\textsuperscript{(13)} al unui pădurar râma cu lăcomie, doar-doar le va găsi.

%%K t09-tit-10 sub
\VUsaut{3.0mm}
\VUpk{10.2pt}{40}{Braconajul.}
\VUsaut{3.0mm}

%%K t09-c2-05-1 p
5. --- Am prins \VUgras{sfârșitul} unei scene care m-a înveselit mult. Datorită
mirosului \VUgras{teckelului} său \textsuperscript{(15)}, \VUgras{pădurarul}
\textsuperscript{(16)} tocmai îl prinsese pe \VUgras{braconier}
\textsuperscript{(22)} în clipa când acesta se pregătea să ducă \VUgras{vânatul}
\textsuperscript{(17)} ucis de el. Voind mai bine să se despartă de pradă decât să
fie prins, braconierul a fugit, iar \VUgras{iepurele} \textsuperscript{(18)},
\VUgras{cerboaica} \textsuperscript{(19)}, \VUgras{căprioara}
\textsuperscript{(20)} și \VUgras{mistrețul} \textsuperscript{(21)} au rămas culcați
în iarbă, spre marea bucurie a pădurarului.

%%K t09-c2-06-1 p
6. --- Uimește felurimea copacilor din această pădure. \VUgras{Fagii}
\textsuperscript{(23)} și \VUgras{mestecenii} \textsuperscript{(26)},
\VUgras{pinii}, care dau \VUgras{rășină} \textsuperscript{(27)},
\VUgras{stejarii} \textsuperscript{(29)} și mulți alții își împletesc crengile.

%%K t09-c2-06-2 p
\VUgras{Iedera} \textsuperscript{(24)}, care se agață de trunchiurile copacilor
înalți, dă \VUgras{pădurii înalte} \textsuperscript{(25)} un verde viu, care se
îmbină plăcut cu tonul palid și moale al \VUgras{mușchiului}
\textsuperscript{(31)}, unde \VUgras{ghionoaia verde} \textsuperscript{(30)} caută
insecte.

%%K t09-tit-11 sub
\VUsaut{3.0mm}
\VUpk{10.2pt}{40}{Priveliște de pădure.}
\VUsaut{3.0mm}

%%K t09-c2-07-1 p
7. --- Stejarii bătrâni m-au fermecat. \VUgras{Rădăcinile} lor mari și noduroase
\textsuperscript{(32)}, pe jumătate ieșite din pământ, le dau o înfățișare minunată
de putere și de tărie, și nu înțeleg cum se poate hotărî \VUgras{stăpânul}
\textsuperscript{(35)} să îi doboare și să îi dea \VUgras{fierăstrăului}
\textsuperscript{(34)} \VUgras{tăietorului de lemne} \textsuperscript{(33)}.
Bietele \VUgras{trunchiuri} \textsuperscript{(36)}, lipsite de \VUgras{scoarță}
\textsuperscript{(37)} sau despicate cu \VUgras{securea} \textsuperscript{(38)}
\VUgras{zilierului} \textsuperscript{(39)}, mă întristează.

%%K t09-c2-08-1 p
8. --- Alături un \VUgras{cărbunar} bătrân \textsuperscript{(41)} a făcut cu
\VUgras{lopata} sa o \VUgras{grămadă} \textsuperscript{(42)} de cărbune, iar o
bătrână ducea o \VUgras{legătură} \textsuperscript{(40)} de \VUgras{vreascuri} din
crengile copacilor tăiați.

%%K t09-tit-12 sub
\VUsaut{3.0mm}
\VUpk{10.2pt}{40}{Vânătoarea.}
\VUsaut{3.0mm}

%%K t09-c2-09-1 p
9. --- Aveam de gând să mă odihnesc puțin în \VUgras{casa pădurarului}
\textsuperscript{(46)}, dar, ajungând la \VUgras{iazul} mic \textsuperscript{(43)},
pe care înota o \VUgras{lebădă} frumoasă \textsuperscript{(45)}, am băgat de seamă
că \VUgras{revărsarea} de curând \textsuperscript{(44)} prefăcuse poteca într-o
adevărată \VUgras{mlaștină}. A trebuit să fac un ocol, și, trecând pe lângă
\VUgras{cedrul} \textsuperscript{(49)}, \VUgras{plopii} \textsuperscript{(48)} și
\VUgras{ulmul} \textsuperscript{(47)} care stau la marginea pădurii, am văzut cum
din \VUgras{desiș} \textsuperscript{(54)} a țâșnit un \VUgras{cerb} minunat
\textsuperscript{(50)}, urmărit de o \VUgras{haită} neobosită
\textsuperscript{(51)}, pe care o înfierbânta \VUgras{cornul}
\textsuperscript{(53)} \VUgras{vânătorilor călări} \textsuperscript{(52)}. Bietul
animal era cu totul sleit. A căzut, dându-și sufletul, la câțiva pași de mine.

%%K t09-c2-10-1 p
10. --- Fiul pădurarului, atras de zgomotul \VUgras{vânătorii}, a ieșit din casă,
și ne-am întors amândoi în pădure. A zărit o \VUgras{coțofană}
\textsuperscript{(56)} pe creanga unui stejar bătrân. Credea că cuibul ei este
ascuns în \VUgras{frunziș} \textsuperscript{(58)} și voia să îl ia.

%%K t09-c2-10-2 p
Sprinten ca o \VUgras{veveriță} \textsuperscript{(61)}, se urca din \VUgras{creangă}
\textsuperscript{(55)} în creangă, dar nu a găsit nimic și a speriat numai o
\VUgras{cioară} bătrână \textsuperscript{(60)}, care ședea pe o \VUgras{cracă}
\textsuperscript{(57)}.
\VUsaut{4.0mm}
\VUfilet{22mm}
